% Local proofs for the explicitly typed Novaes manuscript reconstruction.
% They do not certify every historical original or the final2008 version.
\begin{annexitem}{annex:p13.latin.descents}
 {thm:p13-latin-finite-descents}{Finite expansions, witness scope, and failures of converse descent}
\begin{proof}
Let \(S=\{s_1,\ldots,s_m\}\) and \(P=\{p_1,\ldots,p_k\}\),
with \(m,k\geq1\), in a classical first-order structure with equality.
Use predicate letters \(S,P\) for these extensions.
A universal \(\forall x(Sx\to Px)\) holds exactly when every \(s_i\)
belongs to \(P\), equivalently when \(\bigwedge_iP(s_i)\) holds.
An existential \(\exists x(Sx\land Px)\) holds exactly when at least
one \(s_i\) belongs to \(P\), equivalently \(\bigvee_iP(s_i)\).
The same intersection is nonempty exactly when some \(p_j\) belongs
to \(S\), giving \(\bigvee_jS(p_j)\).

The negative universal \(\forall x(Sx\to\neg Px)\) holds exactly when
\(S\cap P=\varnothing\), equivalently \(\bigwedge_i\neg P(s_i)\).
Symmetry of intersection gives the equivalent predicate-side expansion
\(\bigwedge_j\neg S(p_j)\).
Because the enumeration of \(P\) is exhaustive, \(P(s_i)\) is equivalent
to \(\bigvee_j(s_i=p_j)\). Thus the affirmative universal also has
the expansion \(\bigwedge_i\bigvee_j(s_i=p_j)\): the selected \(j\)
may depend on \(i\).

If \(\exists x(Sx\land\neg Px)\), choose its witness \(a\in S\setminus P\).
For every \(p_j\), \(a\ne p_j\), so the same witness verifies each
\(\exists x(Sx\land x\ne p_j)\). This proves the claimed conjunctive
negative descent. The converse fails when \(S=P=\{a,b\}\) and \(a\ne b\):
for \(p_j=a\) choose \(b\), and for \(p_j=b\) choose \(a\).
Each conjunct holds, while no member of \(S\) lies outside \(P\).
The varying witnesses explain the failure.

In this same two-object example, the affirmative universal is true,
but no single \(p_j\) equals both \(a\) and \(b\).
Therefore \(\bigvee_j\bigwedge_i(s_i=p_j)\) is false.
Moving the predicate alternatives outside all subject predications
strengthens the formula and is not licensed by the original universal.
These proofs use finite exhaustive extensions and classical equality;
they assert no historical completeness theorem for supposition theory.
\end{proof}
\end{annexitem}

\begin{annexitem}{annex:p13.latin.collective}
 {thm:p13-latin-collective}{Collective cardinality separates the nominal patterns}
\begin{proof}
Take twelve distinct objects \(a_1,\ldots,a_{12}\), their collection
\(C=\{a_1,\ldots,a_{12}\}\), and the collective predicate
\(N(X)\) defined by \(|X|=12\) for finite \(X\).
Then \(N(C)\) is true. For each \(i\), the singleton \(\{a_i\}\)
has cardinality one, hence \(N(\{a_i\})\) is false.
Their disjunction is consequently false as well.
Thus the collective predication is true while every singular predication
and their disjunction fail. Ordinary propositional conjunction of
twelve formulas would entail a disjunction of those formulas; this
example is not such a conjunction. It is one predication of a collection.
The distinction establishes the stated limitation without positing
a complete historical collective-supposition semantics.
\end{proof}
\end{annexitem}

\begin{annexitem}{annex:p13.latin.negation-context}
 {thm:p13-latin-negation-context}{The two required output modes}
\begin{proof}
Suppose an output function \(f\) consults only the old mode of a term.
The predicate \(B\) in the unnegated negative universal
has distributive supposition. Its contradictory recovers the particular
affirmative, where the stated classification requires determinate
supposition. Thus \(f(\mathrm{Dist})=\mathrm{Det}\).
The predicate \(B\) in the unnegated particular negative also has
distributive supposition. Its contradictory recovers the affirmative
universal, where the stated classification requires merely-confused
supposition. Thus \(f(\mathrm{Dist})=\mathrm{Conf}\).
If these two output modes differ, both equations cannot hold for a
single-valued \(f\). Hence no mode-only rule has the demanded effect.
The conclusion is conditional on the explicitly specified classifications,
not on an independently established identity of all original theories.
Consulting the other term or transforming the whole syntactic context
supplies information excluded from \(f\)'s hypothesis.
\end{proof}
\end{annexitem}

\begin{annexitem}{annex:p13.latin.consequence-operations}
 {thm:p13-latin-consequence-operations}{Every displayed consequence operation and both negative propagations}
\begin{proof}
All evaluations are in the fixed class \(W\) declared in the body.
For transitivity, if a valuation makes \(A\) true, the first consequence
makes \(B\) true and the second makes \(C\) true.
For derived cut 2-prime, a valuation satisfying \(A\) also satisfies \(B\);
it therefore satisfies the joint antecedent \(A,B\) of the other premise,
which gives \(C\).
For derived cut 2-double-prime, a valuation satisfying \(A,C\) satisfies
\(B\) by the first premise and retains \(C\). The joint antecedent
\(B,C\) now gives \(D\). These also prove the intermediate joint-conclusion
steps \(A\Rightarrow A,B\) and \(A,C\Rightarrow B,C\).

No valuation satisfies \(\bot\), so there is no countervaluation to
\(\bot\Rightarrow A\). Every valuation satisfies \(T\), so there is
no countervaluation to \(A\Rightarrow T\).
For concession propagation, if all valuations satisfy \(A\), the
consequence \(A\Rightarrow B\) makes all satisfy \(B\).
For denial propagation, if all satisfy \(\neg B\), any valuation
satisfying \(A\) would satisfy \(B\), a contradiction. Thus all satisfy
\(\neg A\).

For contraposition, a valuation satisfying \(\neg B\) cannot satisfy \(A\),
since \(A\Rightarrow B\). Hence it satisfies \(\neg A\).
If also \(\neg A\Rightarrow C\), it satisfies \(C\), proving
\(\neg B\Rightarrow C\).
If instead \(C\Rightarrow\neg B\), any valuation satisfying \(C\)
satisfies \(\neg B\), then \(\neg A\), proving \(C\Rightarrow\neg A\).
Conjunction projection follows because a true conjunction has both
conjuncts true. Disjunction introduction follows because a true \(A\)
makes \(A\lor B\) true. For conditional detachment, a valuation satisfying
both \(A\to B\) and \(A\) has \(B\) true by the classical conditional
truth condition.

For the first English non-consequence principle, suppose
\(A\Rightarrow B\) and \(A\nRightarrow C\).
The latter supplies \(v\in W\) with \(A\) true and \(C\) false.
The first premise makes \(B\) true at \(v\). Thus \(v\) witnesses
\(B\nRightarrow C\).
For the second, suppose \(A\Rightarrow B\) and \(C\nRightarrow B\).
Take \(v\in W\) making \(C\) true and \(B\) false. It cannot make
\(A\) true, so it witnesses \(C\nRightarrow A\).
These are distinct countermodel transports.

All arguments remain valid when \(W\) is explicitly restricted to a
temporal/background case. A proposition false or true throughout that
restricted class supplies its relative impossible or necessary boundary.
Nothing in these proofs makes it false or true throughout a larger class.
No proof of completeness of a medieval rule list is asserted.
\end{proof}
\end{annexitem}

\begin{annexitem}{annex:p13.latin.superior}
 {thm:p13-latin-superior}{The included term and its existential condition}
\begin{proof}
If \(x\in M\), inclusion in \(A\) gives \(x\in A\), and inclusion of
\(A\) in \(R\) gives \(x\in R\). Thus \(M\subseteq R\).
If \(M\ne\varnothing\), choose \(m\in M\). The same argument gives
\(m\in R\), proving existence of a member in both.
For the converse counterexample take distinct \(a,b\),
\(M=\{a\}\), \(A=\{a,b\}\), and \(R=\{a\}\).
Every \(M\) runs, and an \(M\) runs, but \(b\in A\setminus R\).
Thus neither conclusion licenses universal predication over \(A\).
If \(M=\varnothing\), the universal inclusion still holds but its
existential conclusion fails, establishing why the import condition
must remain explicit.
\end{proof}
\end{annexitem}

\begin{annexitem}{annex:p13.latin.explosion}
 {thm:p13-latin-explosion}{Every step of the Lewis-style dependency}
\begin{proof}
Start with \(P\land\neg P\). Conjunction projection yields \(P\),
and its other projection yields \(\neg P\).
Disjunction introduction from \(P\) yields \(P\lor Q\), for arbitrary \(Q\).
Disjunctive syllogism from \(P\lor Q\) and \(\neg P\) yields \(Q\).
Every premise of every step is supplied by the preceding steps.
Therefore a calculus admitting these projections, this introduction,
this syllogism and their unrestricted composition has the displayed
derivation for every \(Q\).
Refusing the final consequence while retaining that very derivation
would be incoherent; at least one operation or its scope must be changed.
This is a dependency argument, not a proposed replacement medieval calculus.
\end{proof}
\end{annexitem}

\begin{annexitem}{annex:p13.latin.burley-consistency}
 {thm:p13-latin-burley-consistency}{Admissibility, response uniqueness, and finite consistency preservation}
\begin{proof}
Semantic following is classical. A consistent base means a base with
at least one model. If a base \(\Gamma\) entails both \(\phi\) and
\(\neg\phi\), any model of \(\Gamma\) would satisfy a contradiction,
so no consistent base does so.
Consequently the two pertinent response cases are exclusive.
When neither follows, consistency of \(K_C\) likewise prevents both
knowledge responses from being required. Either exactly one is required
or neither is, in which latter case doubt is required. This proves
response uniqueness at a consistent reached state.

The source's extension fact can be proved in both directions directly.
If \(\Gamma\cup\{\phi\}\) has no model, every model of \(\Gamma\)
must satisfy \(\neg\phi\), so \(\Gamma\Vdash\neg\phi\).
If \(\Gamma\cup\{\neg\phi\}\) also has no model, every model of
\(\Gamma\) must satisfy \(\phi\). Both consequences would contradict
consistency of \(\Gamma\).
Thus at least one extension is consistent.
More particularly, if neither \(\phi\) nor \(\neg\phi\) follows from
\(\Gamma\), there is a model of \(\Gamma\) falsifying each.
One therefore verifies \(\Gamma\cup\{\neg\phi\}\), and the other
verifies \(\Gamma\cup\{\phi\}\). Both extensions are consistent.

The admitted initial positum is satisfiable, so \(\Gamma_0=\{\phi_0\}\)
is consistent. Assume \(\Gamma_{n-1}\) consistent.
If a positive consequence is conceded, all its old models remain models
after the addition. If a negative consequence is denied, its contradictory
literal is added, again preserving all old models.
If the proposal is independent, either knowledge-selected literal
has a model by the preceding paragraph. If doubt is selected, the
base is unchanged. Each branch therefore preserves consistency.
Induction gives consistency after every finite number of correct replies.
By the stipulated outcome condition, Opponent cannot establish a
contradictory reached base. This is the logical response policy claimed.
It gives no time bound for recognizing all semantic dependencies.
\end{proof}
\end{annexitem}

\begin{annexitem}{annex:p13.latin.rome}
 {thm:p13-latin-rome}{All five rounds and the latent contradiction}
\begin{proof}
Use first-order equality and the unary predicate \(L\).
Initially \(\Gamma_0=\{L(a)\lor a=b\}\).
It neither entails \(L(a)\) nor its negation:
a model with \(L(a)\) true verifies the first possibility, while
a model with \(a=b\) and \(L(a)\) false verifies the second.
The knowledge-selected denial therefore gives
\(\Gamma_1=\{L(a)\lor a=b,\neg L(a)\}\).
Disjunctive syllogism yields \(a=b\).
Equality substitution gives \(L(a)\leftrightarrow L(b)\);
with \(\neg L(a)\) this yields \(\neg L(b)\).
Thus \(L(b)\) is a negatively pertinent third proposal and must be denied.

That denial adds only an already entailed literal.
The fourth proposal \(a=b\) is already positively pertinent and is conceded.
The fifth proposal \(L(a)\) remains negatively pertinent and is denied.
Every reached commitment is satisfied by the one-object interpretation
in which \(a,b\) co-denote its single element and \(L\) is empty.
This proves satisfiability of the whole correct trace.

If the third proposal is instead conceded, the base contains
\(a=b,\neg L(a),L(b)\). Equality substitution gives \(L(a)\),
contradicting its negation. Inconsistency is therefore already latent
at the third mistaken reply; the fifth proposal displays it explicitly.
The bad classification of the third proposal, not the unavoidable
correct classification of the fifth, creates the failure.
The shared information \(\neg L(a),L(b)\) has its own model with distinct
denotations of \(a,b\); its consistency is not at issue.
It is not globally united with the counterfactual commitment base.
\end{proof}
\end{annexitem}

\begin{annexitem}{annex:p13.latin.model-restriction}
 {thm:p13-latin-model-restriction}{The two inclusions, properness, and nonemptiness}
\begin{proof}
For an added literal \(\psi\), a model satisfies
\(\Gamma_n\cup\{\psi\}\) exactly when it satisfies both \(\Gamma_n\)
and \(\psi\). This proves the intersection equation.
In particular every new model is an old model.
If \(\Gamma_n\Vdash\psi\), every old model satisfies \(\psi\)
and hence the new base. This proves the reverse inclusion and equality,
including concession of a positive consequence and denial of a negative one.

If \(\phi\) was independent and is conceded, non-entailment supplies
an old model falsifying \(\phi\), which is excluded from the new class.
Independence from \(\neg\phi\) supplies an old model satisfying \(\phi\),
which remains. Hence the inclusion is proper and the new class nonempty.
If independent \(\phi\) is denied, interchange its positive and negative
literals: one old model verifies \(\phi\) and is removed, and another
verifies \(\neg\phi\) and remains. Thus the same conclusion holds.
Doubt adds no literal, so the classes are identical.
These prove both source results and all response branches, without
equating two classes merely because both satisfy the same displayed base.
\end{proof}
\end{annexitem}

\begin{annexitem}{annex:p13.latin.swyneshed-response}
 {thm:p13-latin-swyneshed-response}{The partition and state-dependent function}
\begin{proof}
For each formula \(\phi\), either at least one of
\(\phi_0\Vdash\phi,\phi_0\Vdash\neg\phi\) holds or neither holds.
These alternatives are respectively membership in \(\Delta_{\phi_0}\)
and \(\Pi_{\phi_0}\), proving exhaustiveness and disjointness.
Satisfiability of \(\phi_0\) prevents both pertinent entailments.
Thus a pertinent formula has one prescribed answer, apart from the
separately stipulated pragmatic re-positing exception.
For an impertinent formula, consistency of \(S_n\) prevents both
knowledge entailments. If one holds, its answer is prescribed;
if neither holds, the answer is doubt. Hence the ordinary response
relation is a single-valued total function.

The conditions mention only \(\phi_0,\phi,S_n\), not the previous log.
Two occurrences with those same three data therefore receive the same
reply, irrespective of other proposals between them.
For possible state dependence take independent atoms \(r,p\),
positum \(r\), and knowledge states \(S_1=\{p\}\), \(S_2=\{\neg p\}\).
The positum entails neither \(p\) nor \(\neg p\), so \(p\) is impertinent.
It is conceded in the first state and denied in the second.
This verifies the qualification about changing states.

Finally, contingency cannot be tested by a uniform finite reply record.
The formula \(p\) is contingent over ordinary valuations, since there
are valuations making it true and false. Yet in a sequence of states
all containing \(p\), every external knowledge reply about \(p\) is
positive; and as an admitted positum its ordinary repeated reply is
positive as well. Absence of observed reply variation therefore does
not establish necessity. The reply-history proposal is not used as
the operative modal-admissibility test.
\end{proof}
\end{annexitem}

\begin{annexitem}{annex:p13.latin.swyneshed-connectives}
 {thm:p13-latin-swyneshed-connectives}{Both full reply computations}
\begin{proof}
First let \(\phi_0=\neg(p\land q)\).
Valuations \(p=1,q=0\) and \(p=0,q=1\) satisfy the positum, but
disagree on each individual atom. Neither atom nor its contradictory
is entailed; both atoms are impertinent.
The knowledge state \(\{p,q\}\) is consistent and entails both,
so each receives reply 1.
The conjunction is repugnant to the positum and receives reply 0.
The positum is admissibly contingent: valuation \(0,0\) satisfies it
and valuation \(1,1\) falsifies it.
The combined commitments contain \(p,q,\neg(p\land q)\) and have
no classical model.

Second let \(\phi_0=p\lor q\).
The same valuations \(1,0\) and \(0,1\) show that \(p,q\) are individually
impertinent. The consistent state \(\{\neg p,\neg q\}\) requires denial
of both, giving replies 0,0. Their disjunction is the admitted positum,
not a pragmatically exceptional one, and is conceded with reply 1.
The positum is contingent because \(1,0\) verifies it and \(0,0\) does not.
The combined commitments \(p\lor q,\neg p,\neg q\) have no model.

Within each evaluation the ordinary truth tables were used throughout.
The discrepant replies arise from switching between the positum and
knowledge state, not from a valuation making a true conjunction have
a false conjunct or a true disjunction have both disjuncts false.
\end{proof}
\end{annexitem}

\begin{annexitem}{annex:p13.latin.partial-knowledge}
 {thm:p13-latin-partial-knowledge}{A same-column counterexample under incomplete information}
\begin{proof}
Take distinct independent atoms \(r,p,q\), positum \(r\), and
\(S=\operatorname{Cn}(\{p\lor q\})\), classical consequence closure.
Models with \(r=1,p=q=0\) and \(r=1,p=1,q=0\) show that the positum
entails neither the disjunction nor its negation.
Models with \(r=1,p=1,q=0\) and \(r=1,p=0,q=1\) show that neither
atom nor its negation is entailed by the positum.
Thus all three formulas are impertinent.

The state \(S\) entails \(p\lor q\) by construction.
A model with \(p=0,q=1\) refutes entailment of \(p\) by \(S\);
one with \(p=1,q=0\) refutes entailment of \(\neg p\).
The analogous valuations refute both entailments about \(q\).
The impertinent response rule therefore gives 1 to the disjunction
and doubt to both disjuncts. The state is consistent because both
displayed mixed valuations are models of it.
Consequently neither consistency nor fixed temporal information implies
that knowing a disjunction supplies knowledge of one disjunct.
A complete valuation state would supply that extra condition, but it
is not part of the declared imperfect-information reconstruction.
\end{proof}
\end{annexitem}
